% =====================================================================
\begin{thebibliography}{99}
\addcontentsline{toc}{chapter}{Bibliografía}
\markboth{Bibliografía}{}
\small

\subsection*{Cursos y problemas universitarios citados}

\bibitem{cs336a1}
Equipo docente de CS336 (T.~Hashimoto y P.~Liang).
\newblock CS336: Language Modeling from Scratch. Tarea~1, «Basics», versión 26.0.3.
\newblock Stanford University, primavera de 2026. Problemas \texttt{transformer\_accounting} y \texttt{adamw\_accounting}.
\newblock \url{https://github.com/stanford-cs336/assignment1-basics}

\bibitem{cs336a2}
Equipo docente de CS336 (T.~Hashimoto y P.~Liang).
\newblock CS336: Language Modeling from Scratch. Tarea~2, «Systems and Parallelism», versión 26.1.3.
\newblock Stanford University, primavera de 2026. §8, problema \texttt{data\_parallel\_calcs}.
\newblock \url{https://github.com/stanford-cs336/assignment2-systems}

\bibitem{cs336a3}
Equipo docente de CS336 (T.~Hashimoto y P.~Liang).
\newblock CS336: Language Modeling from Scratch. Tarea~3, «Scaling Laws», versión 26.0.5.
\newblock Stanford University, primavera de 2026. §2.1, problema \texttt{chinchilla\_isoflops}.
\newblock \url{https://github.com/stanford-cs336/assignment3-scaling}

\bibitem{cs336a4}
Equipo docente de CS336 (T.~Hashimoto y P.~Liang).
\newblock CS336: Language Modeling from Scratch. Tarea~4, «Data: Filtering Language Modeling Data», versión 26.0.1.
\newblock Stanford University, primavera de 2026. Problema \texttt{minhash\_deduplication}.
\newblock \url{https://github.com/stanford-cs336/assignment4-data}

\bibitem{hsym}
J.~Austin, S.~Douglas, R.~Frostig, A.~Levskaya \emph{et al.}
\newblock \emph{How to Scale Your Model: A Systems View of LLMs on TPUs}.
\newblock Google DeepMind, 2025. Parte~1 («All About Rooflines»), problema~5; parte~5 («How to Parallelize a Transformer for Training»), problemas~2 y~3; parte~7 («All About Transformer Inference»), problemas~1, 2, 3 y~5.
\newblock \url{https://jax-ml.github.io/scaling-book/}

\bibitem{mmds}
J.~Leskovec, A.~Rajaraman y J.~D. Ullman.
\newblock \emph{Mining of Massive Datasets}, 3.ª ed.
\newblock Cambridge University Press, 2020. Cap.~3, «Finding Similar Items», ejercicios~3.4.1 y~3.4.2.
\newblock \url{http://infolab.stanford.edu/~ullman/mmds/ch3.pdf}

\bibitem{cs246}
J.~Leskovec.
\newblock CS246: Mining Massive Data Sets.
\newblock Stanford University. Curso cuyo libro de texto es~\cite{mmds}.
\newblock \url{https://web.stanford.edu/class/cs246/}

\subsection*{Artículos y fichas técnicas}

\bibitem{hoffmann}
J.~Hoffmann, S.~Borgeaud, A.~Mensch \emph{et al.}
\newblock Training compute-optimal large language models.
\newblock En \emph{Advances in Neural Information Processing Systems 35 (NeurIPS)}, 2022.
\newblock \url{https://arxiv.org/abs/2203.15556}

\bibitem{kaplan}
J.~Kaplan, S.~McCandlish, T.~Henighan \emph{et al.}
\newblock Scaling laws for neural language models.
\newblock Prepublicación, 2020.
\newblock \url{https://arxiv.org/abs/2001.08361}

\bibitem{llama3}
A.~Grattafiori \emph{et al.} (Llama Team, AI @ Meta).
\newblock The Llama 3 herd of models.
\newblock Informe técnico, 2024.
\newblock \url{https://arxiv.org/abs/2407.21783}

\bibitem{schaeffer}
R.~Schaeffer, B.~Miranda y S.~Koyejo.
\newblock Are emergent abilities of large language models a mirage?
\newblock En \emph{Advances in Neural Information Processing Systems 36 (NeurIPS)}, 2023.
\newblock \url{https://arxiv.org/abs/2304.15004}

\bibitem{palm}
A.~Chowdhery, S.~Narang, J.~Devlin \emph{et al.}
\newblock PaLM: Scaling language modeling with Pathways.
\newblock \emph{Journal of Machine Learning Research}, 24, 2023.
\newblock \url{https://arxiv.org/abs/2204.02311}

\bibitem{zero}
S.~Rajbhandari, J.~Rasley, O.~Ruwase e Y.~He.
\newblock ZeRO: Memory optimizations toward training trillion parameter models.
\newblock En \emph{Proceedings of SC20}, 2020.
\newblock \url{https://arxiv.org/abs/1910.02054}

\bibitem{flash}
T.~Dao, D.~Y. Fu, S.~Ermon, A.~Rudra y C.~Ré.
\newblock FlashAttention: Fast and memory-efficient exact attention with IO-awareness.
\newblock En \emph{Advances in Neural Information Processing Systems 35 (NeurIPS)}, 2022.
\newblock \url{https://arxiv.org/abs/2205.14135}

\bibitem{grok1}
xAI.
\newblock Open release of Grok-1, 17 de marzo de 2024.
\newblock \url{https://x.ai/news/grok-os}

\bibitem{leviathan}
Y.~Leviathan, M.~Kalman e Y.~Matias.
\newblock Fast inference from transformers via speculative decoding.
\newblock En \emph{Proceedings of the 40th International Conference on Machine Learning (ICML)}, 2023.
\newblock \url{https://arxiv.org/abs/2211.17192}

\bibitem{vllm}
W.~Kwon, Z.~Li, S.~Zhuang \emph{et al.}
\newblock Efficient memory management for large language model serving with PagedAttention.
\newblock En \emph{Proceedings of the 29th Symposium on Operating Systems Principles (SOSP)}, 2023.
\newblock \url{https://arxiv.org/abs/2309.06180}

\subsection*{Apuntes del curso}

\bibitem{apuntes}
Anatomía de un LLM. Apuntes de los módulos 16 («Preentrenamiento a escala»), 17 («Hardware y sistemas»), 18 («Arquitecturas modernas») y 19 («Inferencia»), edición de octubre de 2026.

\end{thebibliography}
